\section{综合讨论}\label{sec:discussion}

\subsection{两条线的对照}

\begin{table}[H]
\centering
\caption{同一组问题在两个任务上的形态（各行的图集/数据不同，见对应章节）}
\label{tab:both}
\small
\setlength{\tabcolsep}{4pt}
\begin{tabular}{p{2.6cm}p{5.2cm}p{5.2cm}}
\toprule
 & 第一部分：最短路 & 第二部分：节点分类 \\
\midrule
白盒是否存在 & 是（Dijkstra / Bellman-Ford） & 否 \\
三轴怎么量 & 与精确解直接比 & 只能比结构；解释的真值靠合成环境 \\
准确率 & 白盒 $\sim10^{-6}$ vs 黑盒 11.77 & GCN 0.8162 vs MLP 0.5539（两者都是黑盒） \\
鲁棒性的定义 & 扰动后误差的\emph{可证}上界 & 攻击后准确率的\emph{下降}（两者不是同一个量） \\
鲁棒性的结果 & 三条命题：删边单调/非树边不变各 500/500、上界 398 例零违反 & FGA 掉 11.9 点 vs 随机掉 1.7 点 \\
可解释性 & 路径即计算本身，成路率 100\% & 合成环境 0.939 效率；真实数据只能测充分/必要/稳定 \\
本部分的结论 & 预算充足时白盒占优；“学习”是负收益 & 没有白盒，只能退而求其次 \\
\bottomrule
\end{tabular}
\end{table}

表~\ref{tab:both} 是全文的骨架。它想说的一件事是：\textbf{“可解释性 vs. 性能”不是一条普遍的权衡曲线，而是一个条件命题}——曲线只在算法未知时才可能出现。

但要立刻承认一件事：\textbf{本文没有画出这条曲线}。第一部分只证明了“已知时三轴同向、不存在权衡”，第二部分只做到“在没有白盒时先建立起可验证的解释”。要画曲线，需要在一侧扫预算（例如解释稀疏度 $k$ 或攻击预算 $\varepsilon$ 与准确率的 Pareto 前沿），本文没有做。所以这里能说的是一个\emph{必要条件式}的判断：权衡曲线存在的前提是算法未知；已知时它不存在。

\subsection{保真度在两处都失效，但原因不同}

第一部分的反例是：Fidelity$^+$ 0.978（28 图集），而同一集合上的相关系数只有 0.13、更难的 60 图集上是 $-0.10$。失效原因是\textbf{模型没在算}——一个近乎常数的预测器，删掉任何边输出都不变。

这里要补一句第二部分容易被误读的地方：充分性、必要性、稳定性与 Fidelity$^+$ 属于同一族指标（改输入、看输出变多少），第一部分的反例对它们同样成立。第二部分之所以还能用它们，是因为 Cora 上的 GCN 确实在用结构（0.8162 对 MLP 0.5539），不属于“常数预测器”那种退化情形；但这只是\emph{排除了一个已知的反例}，不等于这些指标本身可信。它们是必要条件，不是充分条件。

第二部分的反例是：解释的充分性、必要性、稳定性三项都合格（1.31 vs 1.21、0.123 vs 0.019、Jaccard 0.714），但在 FGA 把准确率打掉 7.5 个点之后，解释\emph{几乎逐条不变}（Jaccard 1.000）。失效原因是\textbf{解释不跟踪模型改变的主意}。

两者共同指向一件事：\textbf{任何单一的解释质量指标都可能被刷}。能做的只有两条——给它配一个“会不会被骗”的对照（第一部分用常数预测器，第二部分用 bottom-$k$），以及尽量找一个有真值的环境（第二部分的 SBM）。

\subsection{可解释性可以怎样被操作化}

本文建议用\textbf{可验证性}替代“可解释性”这个模糊的词：一个解释是否可验证，取决于\emph{能不能被一个独立的、有真值的检验证伪}。

第一部分满足这条：解释边一删除或一放权，真实距离就变，可以逐条核对；白盒的误差还有可计算的上界（命题~\ref{prop:bound}）。
第二部分不满足：Cora 上没有真值，只能造一个 SBM 来测。这不是取巧，而是承认——\textbf{在算法未知的任务上，“解释得对”本身就缺少定义}，我们只能用“在一个人造的可判定环境里解释得对”来近似它。

这也是为什么第二部分的结论比第一部分弱得多：第一部分说的是“白盒更好”，第二部分只能说到“解释器在合成环境里有效，但在真实数据上，稳定不等于忠实”。

\subsection{结论的边界}\label{sec:limits}

这一节的每一条都是对本文结论的削弱，但必须写。

\begin{enumerate}[leftmargin=1.8em]
\item \textbf{第一部分的黑盒基线没有调到最好。}L2 的 6/8 个种子塌缩成常数预测，它最好的一次是 2.19。因此强结论只到“未调好的 mean-pool 基线做不到”，不是“GNN 学不会最短路”。更直接的反例是：Nerem 等人~\cite{nerem2025ood} 证明带合适正则的 GNN 可以\emph{精确}实现 Bellman-Ford 并在规模上外推——这与本文 L0 的结论一致（结构对了就精确），也说明 L2 的失败更像基线问题。
\item \textbf{参数量与迭代预算不对称。}L1 有 2 个参数、跑 128 轮；L2 有 $104{,}321$ 个参数、只有 8 层。这是刻意的（结构对齐的意义就是用先验换样本），但比较因此是联合比较，不是受控消融；表~\ref{tab:same} 已给出同深度的版本，结论明显变弱。
\item \textbf{两个任务都是被刻意选的。}最短路有精确解，是白盒的主场；Cora 是小规模标准图（2708 节点），结论未必推广到大规模异质图。它们的价值在于构成对照，而不是覆盖图学习的全部。
\item \textbf{第二部分的预算是单一的。}只测了 $\varepsilon=5\%$ 的攻击预算与 $k=6$ 的解释预算，没有敏感性曲线；也只训了一个种子做攻击与解释（跨种子稳健性只在解释稳定性一处测过）。
\item \textbf{合成图的“真值”有代价。}SBM 的社区结构与真实图的异质性差距很大。有真值的代价是不真实——这是第二部分所有解释结论的根本限制。
\item \textbf{第一部分的训练只跑了一组超参}（40 轮、lr $10^{-2}$、固定初始化），没有画 $(a,b)$ 的损失面，因此不能排除“换一组设置就能收敛到 1”。我们保留“训练停在 0.9988”这个观察，但不把它上升为“学习必然失败”。
\item \textbf{第二部分被跳过节点的偏差未审计。}45 个候选里有 1 个掩码并列、排不出序，另有 3 个因 oracle 分母为 0 无法定义效率。我们没有记录这些被跳过节点的度数/候选池分布，因此\emph{无法判断跳过是否有偏}；同时有 4 个节点的本地子图归因与全图归因不一致，被记为候选而没有剔除。这两件事都可能让“93.9\%”略偏乐观。
\item \textbf{实现陷阱曾经让结论反转过。}第二部分的归一化算子写错过两次（度数重复累加、自环写成 $d$ 而非 $d^2$、只传单方向边表），每次都让解释结论反向、而非只是变差。这类“看起来对”的代码只能靠断言与外部复算发现；本文的处理是把这些检查写进 \texttt{check\_node\_gnn.py} 与负控制，但读者仍应把“实现正确性”当作一个独立的风险项。
\end{enumerate}

\subsection{什么任务上黑盒才值得}

按本文的逻辑，黑盒的价值出现在“算法未知或不可写”的地方：目标函数没有多项式算法、输入带有感知噪声、或者人自己也说不清规则只能给例子。
在这些场合，白盒根本没有可写的东西，“误差有没有上界”也就无从谈起——此时黑盒是唯一选择，而可解释性退化为一种\emph{审计}工具：不是去证明它对，而是去发现它在什么情况下会错。

但第二部分的发现给这个期待打了一个折扣：一个在模型明显变差时仍然纹丝不动的解释，做不了审计。要让它有用，解释器必须被设计成\emph{对模型的改变敏感}，而不只是对结构敏感——这是本文没有解决、但值得继续做的问题。
